\section{Elementary normal-tensor and Gram-matrix support}
\label{sec:app_bnt_normal_gram}

Throughout this section, $V^{(N)}(A)$ denotes the length-$N$ matrix product
vector of a tensor $A$, and $O_{AB}(N)$ denotes the corresponding mixed
overlap.  We first record elementary consequences of the two notions of
normality used in Chapter~\ref{ch:bnt}, and then prove the finite-family
linear-independence criteria used in the BNT characterization and matching
arguments.

\begin{lemma}[Positive bond dimensions in a coefficient-form BNT]
    \label{lem:cpsv_bnt_block_dimensions_positive}
    \lean{MPSTensor.IsCPSVBasisOfNormalTensors.blocks_dim_pos}
    \leanok
    \uses{def:bnt_cpsv}
    Every block in a coefficient-form basis of normal tensors has positive
    bond dimension.
\end{lemma}

\begin{proof}\leanok
    \uses{def:bnt_cpsv}
    Eventual linear independence makes the matrix product vector of each block
    nonzero at some positive length. A tensor with zero bond dimension has
    zero matrix product vector at every length, which is a contradiction.
\end{proof}

\begin{lemma}[Bond dimension one implies irreducibility]
    \label{lem:bond_dimension_one_irreducible}
    \lean{MPSTensor.isIrreducibleTensor_of_bondDim_one}
    \leanok
    \uses{def:irreducible_tensor}
    Every tensor with a one-dimensional bond space is irreducible.
\end{lemma}

\begin{proof}\leanok
    The only orthogonal projections on a one-dimensional complex vector space
    are zero and the identity.
\end{proof}

\begin{lemma}[Normality in bond dimension one]
    \label{lem:bond_dimension_one_identity_transfer_normal}
    \lean{MPSTensor.isNormalTensor_of_bondDim_one_of_transferMap_eq_id}
    \leanok
    \uses{def:nt_cpsv, lem:bond_dimension_one_irreducible}
    A tensor with one-dimensional bond space and identity transfer map is a
    normal tensor.
\end{lemma}

\begin{proof}\leanok
    Irreducibility follows from the preceding lemma. The identity transfer map
    has spectral radius one, and its only eigenvalue is one.
\end{proof}

\begin{lemma}[Perron gauges preserve primitivity]
    \label{lem:perron_gauge_primitive_iff}
    \lean{MPSTensor.isPrimitive_transferMap_tpGauge_iff}
    \leanok
    \uses{def:primitive_channel}
    Let $\sigma$ be positive definite. The transfer map of the Perron gauge
    determined by $\sigma$ is primitive if and only if the original transfer
    map is primitive.
\end{lemma}

\begin{proof}\leanok
    The two transfer maps are similar through the invertible congruence
    $X\mapsto\sigma^{-1/2}X\sigma^{-1/2}$. Similar linear maps have the same
    spectrum, so the condition that $1$ is the only eigenvalue on the unit
    circle is unchanged.
\end{proof}

\begin{lemma}[Preservation of an algebraic BNT under positive-length MPV equality]
    \label{lem:bnt_of_same_mpv2_pos}
    \lean{MPSTensor.IsBNT.of_sameMPV₂Pos}
    \leanok
    \uses{def:bnt, def:same_mpv2_pos}
    Let $A$ and $A'$ have the same matrix product vectors at every positive
    length. If $(A_j)_{j=1}^{g}$ is a basis of normal tensors for $A$, then
    the same family is a basis of normal tensors for $A'$.
\end{lemma}

\begin{proof}\leanok
    \uses{def:bnt, def:same_mpv2_pos}
    Substitute the positive-length MPV equality into the spanning identity
    for $A$. Normality and eventual linear independence depend only on the
    family $(A_j)_{j=1}^{g}$.
\end{proof}

\begin{corollary}[Positive-length geometric proportionality phase]
    \label{cor:normal_positive_proportional_mpv_geometric_phase}
    \lean{MPSTensor.NonzeroProportionalMPV₂.%
        exists_unit_phase_power_of_isNormalTensor}
    \leanok
    \uses{def:nt_cpsv, def:mpv_state,
        def:nonzero_proportional_mpv}
    Let $A$ and $B$ be normal tensors of positive, possibly different, bond
    dimensions. Suppose that, for every positive $N$, their matrix product
    vectors are proportional by a nonzero scalar. Then there is $a\in\C$,
    with $|a|=1$, such that
    $\ket{V^{(N)}(A)}=a^N\ket{V^{(N)}(B)}$ for every positive $N$.
    This is the positive-length one-normal-block consequence of
    \cite[Corollary~A.3 and Theorem~2.10]{Cirac2016MPDO_arXiv}.
\end{corollary}

\begin{proof}\leanok
    \uses{lem:nonzero_proportional_mpv_basic,
        thm:normal_proportional_mpv_geometric_phase}
    Positive-length nonzero proportionality implies eventual nonzero
    proportionality. Apply
    Theorem~\ref{thm:normal_proportional_mpv_geometric_phase}.
\end{proof}

\begin{lemma}[Eventual linear independence from finite-index overlap orthonormality]
    \label{lem:bnt_finite_index_overlap_orthonormal_li}
    \lean{MPSTensor.eventually_linearIndependent_of_finite_overlap_tendsto_orthonormal}
    \leanok
    \uses{def:mpv_overlap}
    Let $I$ be a finite index set and let $(A_i)_{i \in I}$ be a finite
    family of tensors such that $O_{A_iA_i}(N) \to 1$ for each $i$, and
    $O_{A_iA_j}(N) \to 0$ whenever $i \neq j$. Then the MPV states
    $\{\ket{V^{(N)}(A_i)}\}_{i \in I}$ are linearly independent for all
    sufficiently large~$N$.
\end{lemma}

\begin{proof}\leanok
    The Gram matrix of the family $\{\ket{V^{(N)}(A_i)}\}_{i \in I}$
    converges entrywise to the identity matrix. Hence for all sufficiently
    large $N$ it is invertible, so the MPV states are linearly independent.
\end{proof}

\begin{theorem}[Existential BNT independence from overlap limits]%
    \label{thm:bnt_li_from_overlap_limits_exists}
    \lean{MPSTensor.exists_eventually_linearIndependent_of_overlap_tendsto_orthonormal}
    \leanok
    \uses{lem:bnt_overlap_orthonormal_li}
    If the self-overlaps of a finite block family tend to $1$ and the
    cross-overlaps of distinct blocks tend to $0$, then the MPV states are
    linearly independent for all sufficiently large system sizes.
\end{theorem}

\begin{proof}\leanok
    \uses{lem:bnt_overlap_orthonormal_li}
    This is Lemma~\ref{lem:bnt_overlap_orthonormal_li}, restated with an
    explicit threshold $N_0$.
\end{proof}

\begin{lemma}[Eventual linear independence for two families]%
    \label{lem:bnt_two_family_overlap_orthonormal_li}
    \lean{MPSTensor.eventually_linearIndependent_of_two_family_overlap_tendsto_orthonormal}
    \leanok
    \uses{def:mpv_overlap, lem:bnt_finite_index_overlap_orthonormal_li}
    Let $(A_j)_{j=1}^{g_A}$ and $(B_k)_{k=1}^{g_B}$ be finite families of
    tensors. Suppose the overlaps inside each family are asymptotically
    orthonormal, and suppose every mixed overlap
    $O_{A_jB_k}(N)$ tends to~$0$. Then the combined MPV family
    \begin{align}
        \{\ket{V^{(N)}(A_j)}\}_{j=1}^{g_A}
        \cup
        \{\ket{V^{(N)}(B_k)}\}_{k=1}^{g_B}
        \notag
    \end{align}
    is linearly independent for all sufficiently large~$N$.
\end{lemma}

\begin{proof}\leanok
    Apply the finite-index Gram-matrix criterion to the disjoint union of the two
    index sets. The within-family hypotheses give the diagonal and
    off-diagonal limits inside each block, and the mixed overlap hypothesis gives
    the off-diagonal limits between the two blocks.
\end{proof}

The following two linear-algebra consequences package the eventual Gram-matrix
criterion and its coefficient-extraction corollary.

\begin{lemma}[Eventual linear independence from overlap orthonormality]
    \label{lem:bnt_overlap_orthonormal_li}
    \lean{MPSTensor.eventually_linearIndependent_of_overlap_tendsto_orthonormal}
    \leanok
    \uses{def:mpv_overlap, lem:bnt_finite_index_overlap_orthonormal_li}
    Let $(A_j)_{j=1}^g$ be a finite family of tensors such that
    $O_{A_jA_j}(N) \to 1$ for each $j$ and $O_{A_iA_j}(N) \to 0$
    for $i \neq j$. Then the MPV states
    $\{\ket{V^{(N)}(A_j)}\}_{j=1}^g$ are linearly independent for all
    sufficiently large~$N$.
\end{lemma}

\begin{proof}\leanok
    Apply Lemma~\ref{lem:bnt_finite_index_overlap_orthonormal_li} with
    $I = \{1,\ldots,g\}$.
\end{proof}

The same conclusion is often more convenient with an explicit threshold rather
than the eventual quantifier.



Eventual linear independence is also what converts an equality of two finite
linear combinations into equality of their coefficients; this is the mechanism
behind the coefficient identities of Chapter~\ref{ch:ft_proof}.

\begin{lemma}[Eventual coefficient extraction from linear independence]%
    \label{lem:eventual_coeff_eq_of_eventual_li}
    \lean{MPSTensor.coefficient_eventually_eq_of_eventually_linearIndependent}
    \leanok
    If a finite family of vectors is linearly independent for all sufficiently
    large $N$, and two finite linear combinations of that family are equal
    for all sufficiently large $N$, then the corresponding coefficients are
    equal for all sufficiently large $N$.
\end{lemma}

\begin{proof}\leanok
    Move the two linear combinations to one side. Linear independence forces
    each coefficient difference to vanish at every sufficiently large length.
\end{proof}

When two BNT families must be compared against each other, the same
Gram-matrix criterion applies to their disjoint union, provided the mixed
overlaps also decay. This is the form used in the exact sector matching of
Section~\ref{sec:sector_bnt_strong_match}.
